\chapter{棱锥与圆锥}\label{ch:g8:solids}

在\cref{ch:g7:areas}中讨论的棱柱与圆柱（截面处处相同的立体图形）之后，
本章转向带顶点的立体图形：棱锥和
圆锥。它们的体积公式含有著名的因子 $\frac13$，本章将说明这一结果的合理性并加以运用；
立体图形的研究在
\cref{ch:g9:solids}中继续，涉及球体与截面。

\section{描述棱锥与圆锥}

\begin{definition}[棱锥、圆锥]\label{def:g8:solids:def}
\begin{itemize}
  \item \emph{棱锥}\index{棱锥}以一个多边形为
        \emph{底面}，以一点为\emph{顶点}，与底面的每个
        顶点分别由三角形面相连。其\emph{高}为
        顶点到底面所在平面的距离。
  \item \emph{圆锥}\index{圆锥}是以圆盘代替多边形，按同样的方式构造所得：
        底圆盘、顶点与弯曲的侧面。
\end{itemize}
若棱锥的底面为正多边形（如正方形、等边三角形、
\dots），且顶点竖直位于
底面中心的正上方，则称该棱锥为\emph{正棱锥}。
\end{definition}

\begin{omfigure}
\begin{tikzpicture}[scale=0.75]
  % square pyramid
  \coordinate (P1) at (-1.3,0);
  \coordinate (P2) at (0.9,0);
  \coordinate (P3) at (1.7,0.75);
  \coordinate (P4) at (-0.5,0.75);
  \coordinate (S) at (0.2,3.0);
  \draw[thick] (P1) -- (P2) -- (P3);
  \draw[dashed] (P3) -- (P4) -- (P1);
  \draw[thick] (P1) -- (S) -- (P2);
  \draw[thick] (P3) -- (S);
  \draw[dashed] (P4) -- (S);
  \draw[dashed] (S) -- (0.2,0.375);
  \node[right, font=\small] at (0.22,1.6) {$h$};
  \node[above, font=\small] at (S) {顶点};
  \node[below, font=\small] at (-0.2,-0.1) {正方形底面};
  % cone
  \begin{scope}[xshift=6.2cm]
  \draw[thick] (0,0) ellipse (1.5 and 0.5);
  \coordinate (T) at (0,3.0);
  \draw[thick] (-1.5,0) -- (T) -- (1.5,0);
  \draw[dashed] (T) -- (0,0);
  \node[right, font=\small] at (0.02,1.5) {$h$};
  \draw[thin] (0,0) -- (1.5,0) node[midway, below, font=\small] {$r$};
  \end{scope}
\end{tikzpicture}

{\small 底面为正方形的棱锥与圆锥：各有一个多边形或圆形底面，
一个顶点，高沿垂直于底面的方向度量。}
\end{omfigure}

\begin{example}[数面、棱与顶点]\label{ex:g8:solids:count}
底面为正方形的棱锥有 $5$ 个面（$1$ 个正方形 $+ 4$
个三角形），$8$ 条棱和 $5$ 个顶点。若底面为 $n$ 边形，则
有 $n + 1$ 个面、$2n$ 条棱、$n + 1$ 个顶点。检验欧拉的一个小规律：
面数 $+$ 顶点数 $=$ 棱数 $+ 2$。
\end{example}

\section{体积}

\begin{theorem}[棱锥与圆锥的体积]\label{thm:g8:solids:volume}
设棱锥或圆锥的底面积为 $B$，高为 $h$，则：
\[
V = \frac{1}{3}\, B \times h
\]
——它是同底、同高的棱柱或圆柱体积的三分之一。
对底面半径为 $r$ 的圆锥，$V = \frac13 \pi r^2 h$。
\end{theorem}
\admitted

\begin{remark}[为什么是三分之一？]
把空心棱锥装满水，再倒入同底、同高的
棱柱中，恰好需要三次才能倒满。更妙的是：把正方体
切成三个全等的棱锥，每个棱锥以正方体的一个面为底、
以正方体的一个顶点为顶点（试着想象一下！）。
一般的证明要用到积分，这是高中分册中的工具。
\end{remark}

\begin{example}\label{ex:g8:solids:volume}
一个棱锥的底面为边长 $5$ cm 的正方形，高为 $9$ cm：
\begin{enumerate}
  \item 底面积：$B = 5^2 = 25$ cm$^2$；
  \item 体积：$V = \frac13 \times 25 \times 9 = 25 \times 3 = 75$
        cm$^3$。
\end{enumerate}
一个圆锥形冰淇淋甜筒，底面半径 $3$ cm，高 $10$ cm：
\[
V = \frac13 \times \pi \times 3^2 \times 10 = 30\pi \approx 94
\text{ cm}^3 .
\]
\end{example}

\begin{example}[注意区分高]\label{ex:g8:solids:slant}
对圆锥而言，不要把高 $h$（顶点到底面中心、与底面垂直的
线段）与\emph{斜高} $s$（顶点到底圆边缘的
线段）混淆。二者由毕达哥拉斯定理联系起来
（见\cref{thm:g8:pythagoras:direct}），在直角三角形
顶点–底面中心–底圆边缘中，有
\[
s^2 = h^2 + r^2 .
\]
因此，$r = 3$、$s = 5$ 的圆锥，其高
为 $h = \sqrt{25 - 9} = 4$，体积为
$\frac13 \pi \times 9 \times 4 = 12\pi$。
\end{example}

\begin{method}[体积问题]\label{met:g8:solids:method}
\begin{enumerate}
  \item 识别立体图形（棱柱、圆柱：$V = Bh$；棱锥、圆锥：
        $V = \frac13 Bh$）；
  \item 先单独计算底面积 $B$；
  \item 检查高与底面垂直（若给出斜高，
        则用毕达哥拉斯定理求高）；
  \item 相乘，保持单位一致，
        必要时最后换算（$1$ L $= 1000$ cm$^3$）。
\end{enumerate}
\end{method}

\section{展开图}

\begin{example}[棱锥的展开图]\label{ex:g8:solids:net}
将正四棱锥展开压平，得到其
\emph{展开图}：正方形底面，四个全等的等腰三角形
分别接在底面各边上。三角形的两条相等边是
棱锥的\emph{侧棱}——它们的长度既不是 $h$，
也不是某个侧面的斜高，故在裁剪硬纸板前须
仔细标注。
\end{example}

\begin{omfigure}
\begin{tikzpicture}[scale=0.62]
  \draw[thick] (0,0) rectangle (3,3);
  \draw[thick] (0,0) -- (1.5,-2.2) -- (3,0);
  \draw[thick] (3,0) -- (5.2,1.5) -- (3,3);
  \draw[thick] (0,3) -- (1.5,5.2) -- (3,3);
  \draw[thick] (0,0) -- (-2.2,1.5) -- (0,3);
  \node[font=\small] at (1.5,1.5) {底面};
\end{tikzpicture}

{\small 正四棱锥的展开图：将四个三角形向上折起，
直至各顶点汇于一点。}
\end{omfigure}

\section{练习}

\begin{exercise}[$\star$]\label{exo:g8:solids:1}
一个底面为三角形的棱锥（即
\emph{四面体}）有多少个面、棱和顶点？底面为六边形时又如何？
\end{exercise}

\begin{exercise}[$\star$]\label{exo:g8:solids:2}
求底面积为 $36$ cm$^2$、高为 $10$ cm 的棱锥的体积，
以及底面为长方形的棱锥的体积，其中长方形底面为
$6 \times 4$ cm、高为 $7.5$ cm。
\end{exercise}

\begin{exercise}[$\star$]\label{exo:g8:solids:3}
求半径为 $5$ cm、高为 $9$ cm 的圆锥的体积
（先用 $\pi$ 表示精确值，再四舍五入到整数 cm$^3$）。
\end{exercise}

\begin{exercise}[$\star$]\label{exo:g8:solids:4}
一个圆柱与一个圆锥的半径均为 $4$ cm，高均为
$12$ cm。求二者的体积，并求它们的体积之比。
\end{exercise}

\begin{exercise}[$\star$]\label{exo:g8:solids:5}
某棱锥的体积为 $56$ cm$^3$，高为 $8$ cm。求其
底面积（写出体积公式并解方程）。
\end{exercise}

\begin{exercise}[$\star$]\label{exo:g8:solids:6}
画出圆锥的展开图（底面为一个圆盘，侧面为一个较大圆盘的
扇形）。圆锥的哪一个量是该扇形的半径：
高 $h$ 还是斜高
$s$？
\end{exercise}

\begin{exercise}[$\star\star$]\label{exo:g8:solids:7}
卢浮宫金字塔的底面为正方形，边长约为 $35$ m，
高约为 $22$ m。估算其体积，精确到最接近的
百立方米。
\end{exercise}

\begin{exercise}[$\star\star$]\label{exo:g8:solids:8}
一个圆锥的斜高为 $13$ cm，底面半径为 $5$ cm。求其
高，再求其体积（用 $\pi$ 表示精确值）。
\end{exercise}

\begin{exercise}[$\star\star$]\label{exo:g8:solids:9}
一个半径为 $4$ cm、高为 $9$ cm 的圆锥形杯，装满至杯口，
然后全部倒入一个半径为 $4$ cm 的圆柱形杯中。
求液体在圆柱形杯中的高度。
\end{exercise}

\begin{exercise}[$\star\star$]\label{exo:g8:solids:10}
一个沙漏由两个全等的圆锥（半径为 $2$ cm，高为
$4.5$ cm）顶点相接构成，全部沙子恰好装满其中一个圆锥。
求沙子的体积（先写精确值，再以 cm$^3$ 为单位，
精确到十分位），以及沙子占沙漏内部总容积的
分数。
\end{exercise}

\begin{exercise}[$\star\star\star$]\label{exo:g8:solids:11}
一个底面为正方形、边长为 $6$ cm 的棱锥，其顶点恰好位于
底面一个顶点的正上方，高为 $8$ cm。
\begin{enumerate}
  \item 公式 $V = \frac13 Bh$ 是否仍然适用？（适用——
        顶点不必居中。求 $V$。）
  \item 求最长侧棱的长度，即由顶点
        到底面上与之相对的顶点的线段（需两次应用毕达哥拉斯定理：先
        求底面对角线）。
\end{enumerate}
\end{exercise}

\section{问题：将正方体切成三个棱锥}

\begin{problem}[{周末问题 --- 因子 $\frac13$
从何而来，以及沿半高切开的棱锥会怎样}]
\label{pb:g8:solids:1}
在\cref{thm:g8:solids:volume}之后的注中断言，一个正方体
可以切成三个全等的棱锥——“试着想象一下！”本题将精确地
描绘这一切分，并由此诚实地得出著名的因子
$\frac13$，然后用第二种方式将正方体切成
\emph{六}个棱锥，最后沿其高的一半处切开一个棱锥——
结果很少有人能一次猜对。

全文中，$ABCDEFGH$ 是棱长为 $a$ 的正方体：底面为 $ABCD$，顶
面 $EFGH$，$E$ 位于 $A$ 上方，$F$ 位于 $B$ 上方，$G$ 位于 $C$ 上方，
且 $H$ 位于 $D$ 上方。

\medskip
\noindent\textbf{第一部分 --- 共顶点的三个棱锥。}
\begin{enumerate}
  \item 顶点 $G$ 属于正方体的三个面。列出
        \emph{不}含 $G$ 的三个面，并
        描述以其中每个面为底、每次都以 $G$ 为顶点所得的三个棱锥，
        其中每个棱锥各有多少个面、
        棱和顶点？
        （见\cref{ex:g8:solids:count}）？
  \item 解释为何将正方体绕其长对角线
        $(AG)$ 旋转三分之一圈后，正方体保持不变，却
        使第 1 题中的三个底面作循环轮换——
        并说明三个棱锥为何因此彼此全等，
        互为完全相同的副本。
  \item 承认三个棱锥互不重叠地充满正方体
        （纸板模型颇具说服力），
        推出每个棱锥的体积。
  \item 将此结果与\cref{thm:g8:solids:volume}中的公式核对：
        底面为 $ABCD$、顶点为 $G$ 的棱锥，其顶点正位于
        底面一个\emph{角}的正上方，
        恰与\cref{exo:g8:solids:11}中的情形一致。
        求出其高，套用公式并加以比较。
  \item 对棱长为 $6$ cm 的正方体，求三个棱锥中每一个的体积，
        并说明这一分解与本章的注水实验
        有何关系
        （即用棱锥注满三次得到一个棱柱）。
\end{enumerate}

\medskip
\noindent\textbf{第二部分 --- 共中心的六个棱锥。}
现设 $O$ 为正方体的中心，将 $O$ 与六个面中每个面的
四个顶点相连。
\begin{enumerate}[resume]
  \item 描述所得的六个棱锥，并说明为何它们
        彼此互为全等的副本。
        每个棱锥的高是多少？
  \item 将正方体的体积在六个棱锥间均分，由此推出每个
        棱锥的体积。
  \item 用公式
        $V = \frac13 B h$ 重新计算该体积，并检验两个结果是否一致。
  \item 两种分解方式对某些非常特殊的棱锥证实了因子 $\frac13$。
        试说明它们为何尚不能
        \emph{证明}对每个棱锥与圆锥都成立的\cref{thm:g8:solids:volume}，
        并说明本章如何处理这一缺口：
        哪一结论被承认，哪一实验对其加以支持，
        以及该结论在丛书哪一册中得到了
        诚实的证明。
  \item 将棱长为 $6$ cm 的正方体切成其六个中心棱锥，
        分别用两种方式计算每个棱锥的体积。
\end{enumerate}

\medskip
\noindent\textbf{第三部分 --- 沿半高切割棱锥。}
$SABCD$ 是正四棱锥：底面边长为 $c$，高为
$h$，顶点 $S$ 位于底面中心的正上方。用平面
截取它，该平面过四条侧棱的\emph{中点}：
$[SA]$、$[SB]$、$[SC]$、$[SD]$。
\begin{enumerate}[resume]
  \item 在三角形 $SAB$ 中应用
        \cref{thm:g8:midpoints:direct}：
        连接 $[SA]$ 与 $[SB]$ 中点的线段，其方向与长度分别是多少？
        由此推出截面是边长为 $\frac{c}{2}$ 的正方形，
        与底面平行，
        且位于高为 $\frac{h}{2}$ 处。
  \item 截面以上的部分本身也是正四棱锥。
        求其底面边长与高，并证明其体积
        恰为原体积 $V$ 的\emph{八分之一}。
  \item 求下半部分（\emph{棱台}，即未完成棱锥的形状）
        的体积。半高处的截面以何种比例
        划分了体积？
  \item 卢浮宫金字塔（见\cref{exo:g8:solids:7}：底面边长
        $35$ m，高 $22$ m）分两阶段清洗：先清洗
        半高以上的玻璃，再清洗半高以下的玻璃。
        半高以下部分占\emph{体积}的几分之几？
        以立方米为单位估算该体积，精确到百位。
  \item 一般性结论：若棱锥的\emph{所有}尺寸都乘以 $2$，
        则其体积会如何变化？
        若每个尺寸都乘以 $k$，体积扩大为原来的多少倍？
        这又如何用一句话解释第 12 题中的
        “八分之一”？
\end{enumerate}
\end{problem}
